\begin{thebibliography}{99}
\bibitem{icao_wp69} IATA, Working Paper 69, MIDANPIRG/21, ICAO. [Online]. Available: \url{https://www.icao.int/Meetings/MIDANPIRG-RASG/Documents/WP\%2069.pdf}
\bibitem{easa_sib} EASA, ``Safety Information Bulletin 2022-02R4: GNSS outages and alterations,'' 3 Jul. 2026.
\bibitem{ukmto} UKMTO, ``Advisory 034-25 and summary report,'' 3--7 Oct. 2025.
\bibitem{mscio} MSCIO, ``Advisory 01/2026, Update 12,'' 2 Apr. 2026.
\bibitem{uae_icao} United Arab Emirates, working paper WP05, 14th Meeting of the MIDANPIRG CNS Sub-Group, ICAO, Oct. 2025.
\bibitem{hesch} J. A. Hesch, D. G. Kottas, S. L. Bowman, and S. I. Roumeliotis, ``Camera-IMU-based localization: Observability analysis and consistency improvement,'' \emph{Int. J. Robot. Res.}, vol. 33, no. 1, pp. 182--201, 2014.
\bibitem{vinsmono} T. Qin, P. Li, and S. Shen, ``VINS-Mono: A robust and versatile monocular visual-inertial state estimator,'' \emph{IEEE Trans. Robot.}, vol. 34, no. 4, 2018.
\bibitem{mer} M. Maimone, Y. Cheng, and L. Matthies, ``Two years of visual odometry on the Mars Exploration Rovers,'' \emph{J. Field Robot.}, vol. 24, no. 3, pp. 169--186, 2007.
\bibitem{madmax} L. Meyer \emph{et al.}, ``The MADMAX data set for visual-inertial rover navigation on Mars,'' \emph{J. Field Robot.}, 2021, doi:10.1002/rob.22016.
\bibitem{ronacher} B. Ronacher, ``Path integration in a three-dimensional world: the case of desert ants,'' \emph{J. Comp. Physiol. A}, vol. 206, no. 3, pp. 379--387, 2020.
\bibitem{wittlinger} M. Wittlinger, R. Wehner, and H. Wolf, ``The ant odometer: stepping on stilts and stumps,'' \emph{Science}, vol. 312, no. 5782, pp. 1965--1967, 2006.
\bibitem{muller88} M. M\"uller and R. Wehner, ``Path integration in desert ants, \emph{Cataglyphis fortis},'' \emph{Proc. Natl. Acad. Sci. USA}, vol. 85, no. 14, pp. 5287--5290, 1988.
\bibitem{sommer04} S. Sommer and R. Wehner, ``The ant's estimation of distance travelled: experiments with desert ants, \emph{Cataglyphis fortis},'' \emph{J. Comp. Physiol. A}, vol. 190, no. 1, pp. 1--6, 2004.
\bibitem{stone} T. Stone \emph{et al.}, ``An anatomically constrained model for path integration in the bee brain,'' \emph{Curr. Biol.}, vol. 27, no. 20, pp. 3069--3085, 2017.
\bibitem{antbot} J. Dupeyroux, J. R. Serres, and S. Viollet, ``AntBot: A six-legged walking robot able to home like desert ants in outdoor environments,'' \emph{Sci. Robot.}, vol. 4, no. 27, eaau0307, 2019.
\bibitem{decroon} G. C. H. E. de Croon, J. J. G. Dupeyroux, S. B. Fuller, and J. A. R. Marshall, ``Insect-inspired AI for autonomous robots,'' \emph{Sci. Robot.}, vol. 7, no. 67, eabl6334, 2022.
\bibitem{wehner03} R. Wehner, ``Desert ant navigation: how miniature brains solve complex tasks,'' \emph{J. Comp. Physiol. A}, vol. 189, no. 8, pp. 579--588, 2003.
\bibitem{muller07} M. M\"uller and R. Wehner, ``Wind and sky as compass cues in desert ant navigation,'' \emph{Naturwissenschaften}, vol. 94, no. 7, pp. 589--594, 2007.
\bibitem{labhart} Labhart Laboratory, ``Mobile robot employing insect strategies for navigation'' (project page). [Online]. Available: \url{https://mls.uzh.ch/en/research/fgrimls/Labhart/ModelsRobots/MobileRobot.html}
\bibitem{lambrinos} D. Lambrinos, R. M\"oller, T. Labhart, R. Pfeifer, and R. Wehner, ``A mobile robot employing insect strategies for navigation,'' \emph{Robot. Auton. Syst.}, vol. 30, no. 1--2, pp. 39--64, 2000.
\bibitem{gattaux} G. G. Gattaux, J. R. Serres, F. Ruffier, and A. Wystrach, ``Visual homing in outdoor robots using mushroom body circuits and learning walks,'' arXiv:2507.09725, 2025.
\bibitem{firlefyn} M. V. M. Firlefyn, J. J. Hagenaars, and G. C. H. E. de Croon, ``Direct learning of home vector direction for insect-inspired robot navigation,'' arXiv:2405.03827, 2024.
\bibitem{horvath} G. Horv\'ath and R. Wehner, ``Skylight polarization as perceived by desert ants and measured by video polarimetry,'' \emph{J. Comp. Physiol. A}, vol. 184, pp. 1--7, 1999.
\bibitem{mourikis} A. I. Mourikis and S. I. Roumeliotis, ``A multi-state constraint Kalman filter for vision-aided inertial navigation,'' in \emph{Proc. IEEE ICRA}, 2007, pp. 3565--3572.
\bibitem{openvins} P. Geneva, K. Eckenhoff, W. Lee, Y. Yang, and G. Huang, ``OpenVINS: A research platform for visual-inertial estimation,'' in \emph{Proc. IEEE ICRA}, 2020, pp. 4666--4672.
\bibitem{orbslam3} C. Campos, R. Elvira, J. J. G\'omez Rodr\'iguez, J. M. M. Montiel, and J. D. Tard\'os, ``ORB-SLAM3: An accurate open-source library for visual, visual-inertial and multi-map SLAM,'' \emph{IEEE Trans. Robot.}, 2021.
\bibitem{huang_fej} G. P. Huang, A. I. Mourikis, and S. I. Roumeliotis, ``A first-estimates Jacobian EKF for improving SLAM consistency,'' in \emph{Experimental Robotics}, Springer, 2009, pp. 373--382.
\bibitem{limourikis} M. Li and A. I. Mourikis, ``High-precision, consistent EKF-based visual-inertial odometry,'' \emph{Int. J. Robot. Res.}, vol. 32, no. 6, pp. 690--711, 2013.
\bibitem{zhang18} Z. Zhang and D. Scaramuzza, ``A tutorial on quantitative trajectory evaluation for visual(-inertial) odometry,'' in \emph{Proc. IEEE/RSJ IROS}, 2018, pp. 7244--7251.
\bibitem{sturm} J. Sturm, N. Engelhard, F. Endres, W. Burgard, and D. Cremers, ``A benchmark for the evaluation of RGB-D SLAM systems,'' in \emph{Proc. IEEE/RSJ IROS}, 2012, pp. 573--580.
\bibitem{geiger} A. Geiger, P. Lenz, and R. Urtasun, ``Are we ready for autonomous driving? The KITTI vision benchmark suite,'' in \emph{Proc. IEEE CVPR}, 2012, pp. 3354--3361.
\bibitem{zhang20} L. Zhang, D. Sidoti, A. Bienkowski, K. R. Pattipati, Y. Bar-Shalom, and D. L. Kleinman, ``On the identification of noise covariances and adaptive Kalman filtering: A new look at a 50 year-old problem,'' \emph{IEEE Access}, vol. 8, pp. 59362--59388, 2020.
\bibitem{amohead} Q. Guo, S. Yang, J. Lyu, J. Sun, D. Lin, and S. He, ``AMO-HEAD: Adaptive MARG-only heading estimation for UAVs under magnetic disturbances,'' arXiv:2510.10979, 2025.
\bibitem{imagenetc} D. Hendrycks and T. Dietterich, ``Benchmarking neural network robustness to common corruptions and perturbations,'' in \emph{Proc. ICLR}, 2019.
\bibitem{failuredrift} A. S. Thomas, S. Agnihotri, and M. Keuper, ``Failure or drift? Evaluating monocular SLAM under synthetic and real-world corruptions,'' arXiv:2608.30690, 2026.
\bibitem{khattak} S. Khattak, C. Papachristos, and K. Alexis, ``Visual-thermal landmarks and inertial fusion for navigation in degraded visual environments,'' arXiv:1903.01656, 2019.
\bibitem{tartanair} W. Wang \emph{et al.}, ``TartanAir: A dataset to push the limits of visual SLAM,'' arXiv:2003.14338, 2020; and TartanAir documentation, \url{https://tartanair.org}.
\bibitem{tartanground} M. Patel \emph{et al.}, ``TartanGround: A large-scale dataset for ground robot perception and navigation,'' arXiv:2505.10696, 2025.
\bibitem{sunbcnn} V. Peretroukhin, L. Clement, and J. Kelly, ``Reducing drift in visual odometry by inferring sun direction using a Bayesian convolutional neural network,'' arXiv:1609.05993, 2016.
\bibitem{joshi} B. Joshi and I. Rekleitis, ``Enhancing visual inertial SLAM with magnetic measurements,'' arXiv:2409.09904, 2024.
\bibitem{zhao} H. Zhao \emph{et al.}, ``Seamless micro-electro-mechanical system-inertial navigation system/polarization compass navigation method with data and model dual-driven approach,'' \emph{Micromachines}, vol. 15, no. 2, p. 237, 2024, doi:10.3390/mi15020237.
\bibitem{nosek} B. A. Nosek, C. R. Ebersole, A. C. DeHaven, and D. T. Mellor, ``The preregistration revolution,'' \emph{Proc. Natl. Acad. Sci. USA}, vol. 115, no. 11, pp. 2600--2606, 2018.
\bibitem{derrac} J. Derrac, S. Garc\'ia, D. Molina, and F. Herrera, ``A practical tutorial on the use of nonparametric statistical tests as a methodology for comparing evolutionary and swarm intelligence algorithms,'' \emph{Swarm Evol. Comput.}, vol. 1, no. 1, pp. 3--18, 2011.
\bibitem{colas} C. Colas, O. Sigaud, and P.-Y. Oudeyer, ``How many random seeds? Statistical power analysis in deep reinforcement learning experiments,'' arXiv:1806.08295, 2018.
\bibitem{suhai} B. Suhai and G. Horv\'ath, ``How well does the Rayleigh model describe the E-vector distribution of skylight in clear and cloudy conditions? A full-sky polarimetric study,'' \emph{J. Opt. Soc. Am. A}, vol. 21, no. 9, pp. 1669--1676, 2004.
\bibitem{pomozi} I. Pomozi, G. Horv\'ath, and R. Wehner, ``How the clear-sky angle of polarization pattern continues underneath clouds: full-sky measurements and implications for animal orientation,'' \emph{J. Exp. Biol.}, vol. 204, pp. 2933--2942, 2001.
\bibitem{liang} H. Liang, H. Bai, T. Zhou, and K. Shen, ``Influence of sensor tilts on bio-inspired polarized skylight orientation determination,'' arXiv:2007.08987, 2020.
\bibitem{noaa} NOAA Global Monitoring Laboratory, ``General solar position calculations.'' [Online]. Available: \url{https://gml.noaa.gov/grad/solcalc/solareqns.PDF}
\bibitem{wang19} Y. Wang, J. Chu, R. Zhang, J. Li, X. Guo, and M. Lin, ``A bio-inspired polarization sensor with high outdoor accuracy and central-symmetry calibration method with integrating sphere,'' \emph{Sensors}, vol. 19, no. 16, 3448, 2019.
\bibitem{potter} B. Potter, E. Mounier, A. Noureldin, Y. Antar, and M. Alam, ``Navigation by skylight polarization via Rayleigh model optimization on a road vehicle,'' arXiv:2610.20797, 2026.
\bibitem{cheng} X. Cheng, G. Liu, A. Hedman, K. Wang, and H. Li, ``Expressway visibility estimation based on image entropy and piecewise stationary time series analysis,'' arXiv:1804.04601, 2018.
\bibitem{tg_hf} TartanGround dataset card, Hugging Face repository theairlabcmu/TartanGround. [Online]. Available: \url{https://huggingface.co/datasets/theairlabcmu/TartanGround}
\bibitem{pust} N. J. Pust and J. A. Shaw, ``Wavelength dependence of the degree of polarization in cloud-free skies: simulations of real environments,'' \emph{Opt. Express}, vol. 20, no. 14, pp. 15559--15568, 2012.
\bibitem{ma15} T. Ma, X. Hu, L. Zhang, J. Lian, X. He, Y. Wang, and Z. Xian, ``An evaluation of skylight polarization patterns for navigation,'' \emph{Sensors}, vol. 15, no. 3, pp. 5895--5913, 2015.
\bibitem{liang20b} H. Liang, H. Bai, and T. Zhou, ``Exploration of whether skylight polarization patterns contain three-dimensional attitude information,'' arXiv:2012.09154, 2020.
\bibitem{hu22} P. Hu, P. Huang, Z. Qiu, J. Yang, and X. Liu, ``A 3D attitude estimation method based on attitude angular partial feedback for polarization-based integrated navigation system,'' \emph{Sensors}, vol. 22, no. 3, 710, 2022.
\bibitem{zhang22} J. Zhang, S. Wang, W. Li, and Z. Qiu, ``A multi-mode switching variational Bayesian adaptive Kalman filter algorithm for the SINS/PNS/GMNS navigation system of pelagic ships,'' \emph{Sensors}, vol. 22, no. 9, 3372, 2022.
\bibitem{sun25} Y. Sun, X. Liu, X. Liu, H. Zhao, C. Wang, H. Cao, and C. Shen, ``Variational Bayesian innovation saturation Kalman filter for MEMS-INS/polarization compass integrated navigation,'' \emph{Micromachines}, vol. 16, no. 9, 1036, 2025.
\bibitem{steidle} F. Steidle, W. St\"urzl, and R. Triebel, ``Visual-inertial sensor fusion with a bio-inspired polarization compass for navigation of MAVs,'' in \emph{Proc. IMAV}, 2019.
\bibitem{chu23} J. Chu, H. Hu, Z. Wan, and J. Li, ``Visual inertial navigation system aided by polarized light,'' \emph{Laser Optoelectron. Prog.}, vol. 60, no. 7, 0726002, 2023.
\bibitem{zeng08} J. Zeng, Q. Han, and J. Wang, ``High spectral resolution simulation of polarization of skylight: Sensitivity to aerosol vertical profile,'' \emph{Geophys. Res. Lett.}, vol. 35, L20801, 2008, doi:10.1029/2008GL035645.
\bibitem{deleon25} C. M. DeLeon, W. R. Espinosa, Y. Long, J. A. Shaw, and M. Kupinski, ``Observing and simulating the polarization neutral point dependence on atmospheric turbidity,'' in \emph{Proc. SPIE 13615, Polarization Science and Remote Sensing XII}, 1361505, 2025, doi:10.1117/12.3062931.
\bibitem{yuan25} M. Yuan, X. Wu, C. Wang, and X. Liu, ``Tilt induced error compensation with vision based method for polarization navigation,'' \emph{Appl. Sci.}, vol. 15, no. 9, 5060, 2025, doi:10.3390/app15095060.
\bibitem{su26} H. Su, S. Zhang, Q. Fu, J. Zhan, Y. Li, C. Wang, and J. Zhou, ``Solar meridian determination based on fused polarization patterns,'' \emph{Chinese Optics}, vol. 19, no. 4, pp. 895--906, 2026, doi:10.37188/CO.2026-0044.
\bibitem{wang24} S. Wang, X. Chi, H. Bai, J. Yang, P. Hu, Z. Qiu, and X. Yu, ``An integrated navigation system using multidirectional polarization sensors and sun sensor for cloudy environments,'' \emph{IEEE Sensors J.}, vol. 24, no. 14, pp. 21906--21915, 2024, doi:10.1109/JSEN.2023.3302405.
\bibitem{sturzl17} W. St\"urzl, ``A lightweight single camera polarization compass with covariance estimation,'' in \emph{Proc. IEEE ICCV}, 2017, doi:10.1109/ICCV.2017.572.
\bibitem{poughon24} L. Poughon, V. Aubry, J. Monnoyer, S. Viollet, and J. R. Serres, ``An extended database of annotated skylight polarization images covering a period of two months,'' \emph{BMC Res. Notes}, vol. 17, 306, 2024, doi:10.1186/s13104-024-06959-6; data set ``1 month long skylight polarization images dataset,'' Recherche Data Gouv, doi:10.57745/610QTS, CC BY-SA 4.0.
\bibitem{giubilato26} R. Giubilato, M. G. M\"uller, M. Sewtz, L. A. Encinar Gonzalez, J. Folkesson, and R. Triebel, ``The S3LI Vulcano dataset: A dataset for multi modal SLAM in unstructured planetary environments,'' arXiv:2601.19557, 2026.
\bibitem{power} NASA Langley Research Center, POWER Project, climatology API v2.10 (aerosol optical depth at 550 nm and cloud amount). [Online]. Available: \url{https://power.larc.nasa.gov}, accessed 3 Nov. 2026.
\end{thebibliography}
